\chapter{Hugging Face and Local Model Deployment}
\label{app:hf}

\begin{remark}[Learning Objectives]
After working through this appendix, the reader should be able to:
\begin{enumerate}
  \item Navigate the Hugging Face Hub to identify models and datasets suited
        to a given financial NLP task.
  \item Load a pre-trained model with the \texttt{pipeline()} API, run
        inference, and interpret the output.
  \item Explain the LoRA and QLoRA fine-tuning strategies, estimate the
        GPU memory required for a given configuration, and apply PEFT to
        adapt a base model to financial text.
  \item Select an appropriate quantization level (GGUF Q4--Q8) for a target
        model size and hardware budget, and deploy the quantised model using
        \texttt{llama.cpp} or Ollama.
  \item Deploy a production inference server with vLLM or TGI and justify
        the choice based on concurrency and latency requirements.
  \item Articulate the data sovereignty, regulatory, and cost arguments for
        on-premise deployment of LLMs in regulated finance environments.
  \item Design a complete local-model pipeline for a quantitative finance
        task---from model selection and optional fine-tuning through to a
        queryable REST endpoint.
\end{enumerate}
\end{remark}

% ---------------------------------------------------------------
\section{The Hugging Face Ecosystem}
\label{app:hf:ecosystem}

\begin{context}
When Meta released LLaMA in February 2023, it uploaded the weights to the
Hugging Face Hub.  Within hours, the model had been downloaded tens of
thousands of times.  That episode crystallised the Hub's role: it had become
the de facto distribution channel for open-weight AI models, the way PyPI is
for Python packages or npm for JavaScript modules.  By 2026 the Hub hosted
more than two million models, half a million datasets, and one million
interactive demo applications.  For finance practitioners, this means that
virtually every open-weight language model relevant to their work---from
general-purpose instruction-following models to domain-specific financial
classifiers---is one \texttt{pip install} and one \texttt{from\_pretrained}
call away.
\end{context}

\subsection{Hub: Models, Datasets, Spaces, and Endpoints}
\label{app:hf:hub}

The Hugging Face Hub (\texttt{huggingface.co}) is a web platform and Python
API that serves four types of artefact:

\begin{description}
  \item[Models] Pre-trained weight checkpoints stored in \texttt{safetensors}
        or \texttt{pytorch\_model.bin} format, versioned with Git LFS.  Each
        model card documents the training data, evaluation results, intended
        use, and known limitations.

  \item[Datasets] Curated tabular, text, vision, and audio datasets accessible
        via the \texttt{datasets} library.  Finance-relevant collections include
        earnings call transcripts, SEC filings, credit ratings histories, and
        financial news corpora.

  \item[Spaces] Hosted Gradio or Streamlit web applications that let model
        authors publish interactive demos without requiring users to run code
        locally.  Useful for evaluating a model's behaviour before committing
        to a download.

  \item[Inference Endpoints] Managed, autoscaling API endpoints for Hub models,
        billed per token.  They provide a production-grade alternative to
        self-hosting for workloads where traffic is variable or the team lacks
        GPU infrastructure.
\end{description}

Interact with the Hub programmatically via the \texttt{huggingface\_hub}
library:

\begin{minted}{python}
from huggingface_hub import hf_hub_download, list_models

# List finance-tagged models sorted by downloads
models = list(list_models(filter="finance", sort="downloads", direction=-1))
for m in models[:5]:
    print(m.id, m.downloads)
\end{minted}

\subsection{Transformers and Datasets Libraries}
\label{app:hf:libraries}

The \textbf{Transformers} library (\texttt{pip install transformers}) is the
primary SDK for working with Hub models.  It provides:

\begin{itemize}
  \item \texttt{AutoTokenizer}, \texttt{AutoModel}, and task-specific heads
        (\texttt{AutoModelForSequenceClassification}, etc.) that load any
        compatible model from the Hub or a local directory.
  \item A \texttt{pipeline()} API that wraps tokenization, inference, and
        post-processing into a single callable.
  \item \texttt{Trainer} and \texttt{SFTTrainer} for supervised fine-tuning.
  \item Tight integration with PEFT, bitsandbytes, and Flash Attention 2.
\end{itemize}

The \textbf{Datasets} library (\texttt{pip install datasets}) provides
memory-mapped access to large datasets via Apache Arrow, enabling streaming
of corpora too large to fit in RAM---an important property when fine-tuning
on multi-year collections of financial filings.

\begin{minted}{python}
from datasets import load_dataset

# Stream SEC 10-K filings without downloading the full corpus
ds = load_dataset("JanosAudran/financial-reports-sec", split="train",
                  streaming=True)
for example in ds.take(3):
    print(example["text"][:200])
\end{minted}

% ---------------------------------------------------------------
\section{Model Discovery and the Pipeline API}
\label{app:hf:models}

\subsection{Browsing the Hub for Finance Tasks}
\label{app:hf:browsing}

The Hub's filter system lets practitioners narrow to models by task, language,
and licence.  For finance NLP the most relevant task tags are:

\begin{description}
  \item[\texttt{text-classification}] Sentiment analysis, intent classification,
        regulatory document categorisation.
  \item[\texttt{token-classification}] Named entity recognition for financial
        entities (companies, currencies, financial instruments).
  \item[\texttt{question-answering}] Extractive QA over earnings reports and
        analyst transcripts.
  \item[\texttt{text-generation}] Instruction-following models for report
        drafting, data extraction, and chain-of-thought reasoning over
        financial data.
  \item[\texttt{feature-extraction}] Embedding models for semantic search over
        document corpora and RAG pipelines.
\end{description}

\subsection{The Pipeline API}
\label{app:hf:pipeline}

The \texttt{pipeline()} function is the quickest path from Hub model to
inference result:

\begin{minted}{python}
from transformers import pipeline

# Zero-shot sentiment on an earnings call excerpt
clf = pipeline("zero-shot-classification",
               model="facebook/bart-large-mnli",
               device=0)          # device=0 → GPU; -1 → CPU

labels = ["positive", "negative", "neutral"]
result = clf(
    "Revenue grew 18% year-over-year, driven by strong mortgage originations.",
    candidate_labels=labels,
)
# {'labels': ['positive', 'neutral', 'negative'], 'scores': [0.89, 0.08, 0.03]}
print(result)
\end{minted}

For generation tasks, set \texttt{max\_new\_tokens} and
\texttt{do\_sample=False} for deterministic output---important in research
contexts where reproducibility matters:

\begin{minted}{python}
gen = pipeline("text-generation", model="mistralai/Mistral-7B-Instruct-v0.3",
               device_map="auto", torch_dtype="auto")

output = gen(
    "<s>[INST] Summarise the key risk factors from: {text} [/INST]",
    max_new_tokens=256, do_sample=False,
)
print(output[0]["generated_text"])
\end{minted}

\subsection{Finance-Specific Model Families}
\label{app:hf:finance-models}

Several model families are specifically relevant to quantitative finance:

\begin{description}
  \item[FinBERT \citep{araci2019finbert}] A BERT model pre-trained on
        financial communication and fine-tuned for financial sentiment
        classification.  Available at
        \texttt{ProsusAI/finbert} on the Hub.  The standard baseline for
        earnings call and news sentiment tasks.

  \item[BloombergGPT \citep{wu2023bloomberggpt}] A 50-billion-parameter model
        trained on Bloomberg's proprietary financial corpus.  Not publicly
        available on the Hub but provides the benchmark against which
        open-weight finance models are evaluated.

  \item[FinLLMs] A family of open-weight models fine-tuned on financial
        instruction datasets; evaluated on FinQA, FPB (Financial Phrase Bank),
        and Headline classification benchmarks.  Multiple variants at different
        parameter scales are available on the Hub.

  \item[Llama 3 / Mistral 7B instruction variants] General-purpose instruction-
        following models that, with appropriate prompting or light fine-tuning,
        perform competitively on financial NLP tasks while remaining locally
        deployable on commodity GPU hardware.
\end{description}

Table~\ref{tab:finance-models} gives a concise evaluation snapshot across
key finance NLP benchmarks.

\begin{table}[ht]
\centering
\caption{Selected model performance on finance NLP benchmarks (higher is better).}
\label{tab:finance-models}
\begin{tabular}{llccc}
\toprule
\textbf{Model} & \textbf{Size} & \textbf{FPB F1} & \textbf{FinQA} & \textbf{Headline} \\
\midrule
FinBERT             & 110M  & 0.88 & ---  & 0.97 \\
Llama 3.1 8B (inst.)& 8B    & 0.82 & 0.61 & 0.93 \\
Mistral 7B (inst.)  & 7B    & 0.80 & 0.58 & 0.91 \\
BloombergGPT        & 50B   & 0.87 & 0.68 & 0.97 \\
\bottomrule
\multicolumn{5}{l}{\footnotesize FPB: Financial Phrase Bank. FinQA: numerical reasoning over financial tables.} \\
\multicolumn{5}{l}{\footnotesize Headline: financial news headline classification. Scores are approximate.}
\end{tabular}
\end{table}

% ---------------------------------------------------------------
\section{Parameter-Efficient Fine-Tuning}
\label{app:hf:peft}

\begin{context}
Full fine-tuning a 7-billion-parameter model requires updating all seven
billion weights simultaneously, which at 16-bit precision demands roughly
14\,GB of GPU memory for the weights alone.  Optimizer states add substantially
more: AdamW stores a float32 master copy of the weights plus first and second
moment tensors, totalling at least 112\,GB for a 7B model in mixed-precision
training.  For most practitioners, this is infeasible.  Parameter-efficient fine-tuning (PEFT) methods solve this by
freezing the pre-trained weights and learning a small number of additional
parameters that steer the model's behaviour.  LoRA reduces the memory
requirement by a factor of 10--20 while retaining 90--95\% of the quality of
full fine-tuning, enabling adaptation on a single consumer GPU.
\end{context}

\subsection{LoRA: Low-Rank Adaptation}
\label{app:hf:lora}

LoRA \citep{hu2022lora} decomposes each weight update $\Delta W \in
\mathbb{R}^{d \times k}$ as the product of two low-rank matrices:
\begin{equation}
  \Delta W = B A, \quad B \in \mathbb{R}^{d \times r},\; A \in \mathbb{R}^{r \times k},\; r \ll \min(d, k).
  \label{eq:lora-hf}
\end{equation}
During training only $A$ and $B$ are updated; the pre-trained weight $W_0$
is frozen.  At inference, the adapter can be merged: $W = W_0 + \Delta W$,
adding zero latency overhead.  The rank $r$ is a hyperparameter: values of
4--16 are typical for fine-tuning general instruction-following models;
higher ranks (32--64) are used when the target domain is very different from
the pre-training distribution.

\begin{minted}{python}
from transformers import AutoModelForCausalLM, AutoTokenizer
from peft import get_peft_model, LoraConfig, TaskType

base = AutoModelForCausalLM.from_pretrained("mistralai/Mistral-7B-v0.3",
                                             torch_dtype="auto",
                                             device_map="auto")

config = LoraConfig(
    task_type=TaskType.CAUSAL_LM,
    r=16,                          # rank
    lora_alpha=32,                 # scaling factor
    target_modules=["q_proj", "v_proj"],
    lora_dropout=0.05,
)

model = get_peft_model(base, config)
model.print_trainable_parameters()
# trainable params: 6,815,744 || all params: 7,248,547,840 || trainable%: 0.09
\end{minted}

\subsection{QLoRA: Quantised Low-Rank Adaptation}
\label{app:hf:qlora}

QLoRA \citep{dettmers2023qlora} combines LoRA with 4-bit NormalFloat (NF4)
quantisation of the frozen base weights.  This reduces the memory footprint
of a 7B model from approximately 14\,GB (bfloat16) to approximately 4--5\,GB,
making fine-tuning feasible on a single consumer GPU with 6--8\,GB of VRAM.

\begin{minted}{python}
from transformers import BitsAndBytesConfig
import torch

bnb_config = BitsAndBytesConfig(
    load_in_4bit=True,
    bnb_4bit_quant_type="nf4",
    bnb_4bit_compute_dtype=torch.bfloat16,
    bnb_4bit_use_double_quant=True,   # nested quantization: saves ~0.4 GB
)

base = AutoModelForCausalLM.from_pretrained(
    "mistralai/Mistral-7B-v0.3",
    quantization_config=bnb_config,
    device_map="auto",
)
\end{minted}

\subsection{Fine-Tuning on Financial Text}
\label{app:hf:finetuning}

A practical fine-tuning pipeline for a financial instruction dataset using the
\texttt{SFTTrainer} from the TRL library:

\begin{minted}{python}
from trl import SFTTrainer
from transformers import TrainingArguments
from datasets import load_dataset

dataset = load_dataset("json",
                       data_files="data/finance_instruct.jsonl",
                       split="train")

args = TrainingArguments(
    output_dir="checkpoints/mistral-finance",
    num_train_epochs=3,
    per_device_train_batch_size=2,
    gradient_accumulation_steps=8,
    learning_rate=2e-4,
    fp16=True,
    logging_steps=10,
    save_strategy="epoch",
)

trainer = SFTTrainer(
    model=model,          # QLoRA model from previous step
    train_dataset=dataset,
    dataset_text_field="text",
    max_seq_length=2048,
    args=args,
)
trainer.train()
\end{minted}

Merging the adapter back into the base weights for deployment:

\begin{minted}{python}
from peft import PeftModel

merged = PeftModel.from_pretrained(base, "checkpoints/mistral-finance/final")
merged = merged.merge_and_unload()
merged.save_pretrained("models/mistral-finance-merged")
\end{minted}

% ---------------------------------------------------------------
\section{Quantization and the GGUF Format}
\label{app:hf:quant}

\subsection{Quantization Levels and Quality Trade-offs}
\label{app:hf:quant-levels}

Post-training quantization compresses model weights from 16-bit floats to
lower-bit integer representations, reducing both memory footprint and
inference latency.  Table~\ref{tab:quant-levels} summarises the commonly
used levels in the GGUF ecosystem.

\begin{table}[ht]
\centering
\caption{GGUF quantization levels: memory and quality trade-offs for a 7B model.}
\label{tab:quant-levels}
\begin{tabular}{lrrl}
\toprule
\textbf{Quantization} & \textbf{VRAM (7B)} & \textbf{Quality loss} & \textbf{Best use} \\
\midrule
Q2\_K     & $\approx$2.8\,GB & High ($>$5\%)  & Ultra-constrained hardware only \\
Q4\_K\_M  & $\approx$4.1\,GB & Low ($<$1\%)   & Standard local deployment \\
Q5\_K\_M  & $\approx$4.8\,GB & Very low       & Accuracy-sensitive tasks \\
Q6\_K     & $\approx$5.5\,GB & Negligible     & Near-lossless on 6\,GB VRAM \\
Q8\_0     & $\approx$7.2\,GB & Negligible     & Research: maximal quality \\
F16       & $\approx$14\,GB  & None (baseline)& Serving; requires server GPU \\
\bottomrule
\end{tabular}
\end{table}

For most finance NLP tasks, Q4\_K\_M delivers the best balance: approximately
75\% memory reduction relative to full precision with less than 1\% quality
degradation on standard benchmarks \citep{llamacpp2023}.

\subsection{GGUF and llama.cpp}
\label{app:hf:llamacpp}

GGUF (GPT-Generated Unified Format) is the standard single-file format for
quantized model checkpoints.  A GGUF file bundles the quantized weights,
the tokenizer vocabulary, positional encoding parameters, and all metadata
needed to run inference---enabling distribution as a single artefact with no
separate configuration files.

\texttt{llama.cpp} is an open-source C/C++ inference engine that reads GGUF
files and runs on CPUs, Apple Silicon (via Metal), and NVIDIA GPUs (via CUDA).
It is the reference implementation for local LLM inference and underpins most
higher-level tools, including Ollama.

\begin{minted}{bash}
# Download a GGUF model from the Hub (bartowski provides many quantizations)
huggingface-cli download bartowski/Mistral-7B-Instruct-v0.3-GGUF \
    --include "Mistral-7B-Instruct-v0.3-Q4_K_M.gguf" \
    --local-dir models/

# Run interactive inference
./llama-cli -m models/Mistral-7B-Instruct-v0.3-Q4_K_M.gguf \
            -n 256 -p "Summarise the credit risk in: {text}"
\end{minted}

\texttt{llama.cpp} also exposes an OpenAI-compatible HTTP server:

\begin{minted}{bash}
./llama-server -m models/Mistral-7B-Instruct-v0.3-Q4_K_M.gguf \
               --port 8080 --n-gpu-layers 35
\end{minted}

\begin{deepdive}
The \texttt{--n-gpu-layers} flag controls how many transformer layers are
offloaded to the GPU; the remainder run on CPU.  For a 7B model with 32
layers on a machine with 4\,GB of GPU VRAM, offloading 20--24 layers
typically maximises throughput while keeping the remainder in system RAM.
The optimal split is model- and hardware-specific; benchmark with
\texttt{llama-bench} before committing to a production configuration.
\end{deepdive}

\subsection{Hardware Requirements by Model Size}
\label{app:hf:hardware}

Table~\ref{tab:hardware-reqs} gives the minimum GPU VRAM needed to run
inference at Q4\_K\_M quantization for the most commonly deployed open-weight
model sizes.

\begin{table}[ht]
\centering
\caption{Minimum GPU VRAM for inference at Q4\_K\_M quantization.}
\label{tab:hardware-reqs}
\begin{tabular}{lrrl}
\toprule
\textbf{Model size} & \textbf{VRAM (Q4\_K\_M)} & \textbf{Example models} & \textbf{Consumer GPU} \\
\midrule
3--4B   & $\approx$2.5\,GB & Phi-3 Mini, Llama 3.2 3B & RTX 3060 (8\,GB) \\
7--8B   & $\approx$4.5\,GB & Mistral 7B, Llama 3.1 8B & RTX 3060 (8\,GB) \\
13B     & $\approx$8\,GB   & Llama 2 13B, CodeLlama 13B & RTX 3080/4070 \\
32--34B & $\approx$20\,GB  & Llama 3.3 32B, CodeLlama 34B & RTX 4090 (24\,GB) \\
70B     & $\approx$42\,GB  & Llama 3.1 70B              & 2$\times$ RTX 4090 or A100 \\
\bottomrule
\end{tabular}
\end{table}

% ---------------------------------------------------------------
\section{Local Deployment with Ollama}
\label{app:hf:ollama}

Ollama provides a user-friendly wrapper around llama.cpp that manages model
downloads, versioning, and a persistent OpenAI-compatible REST server.  By
early 2026 it had accumulated over 112 million pulls for Llama 3.1 alone,
making it the most widely used local LLM runtime in the developer community.

\subsection{Installation and Model Management}
\label{app:hf:ollama-install}

\begin{minted}{bash}
# Install (macOS/Linux)
curl -fsSL https://ollama.com/install.sh | sh

# Pull and run a model (starts the server automatically)
ollama pull mistral               # downloads Mistral 7B Q4_K_M by default
ollama run mistral "Explain duration risk in plain language."

# List locally available models
ollama list
\end{minted}

Ollama stores models in \texttt{\textasciitilde/.ollama/models/} and manages
multiple versions side by side.  Custom GGUF files can be registered with a
\texttt{Modelfile}:

\begin{minted}{text}
# Modelfile for a finance-tuned variant
FROM ./models/mistral-finance-merged.gguf
SYSTEM "You are a quantitative finance assistant. Always use log returns and
        cite data sources."
PARAMETER temperature 0.1
PARAMETER num_ctx 8192
\end{minted}

\begin{minted}{bash}
ollama create mistral-finance -f Modelfile
ollama run mistral-finance "What is the Sharpe ratio of this return series?"
\end{minted}

\subsection{REST API and Python Integration}
\label{app:hf:ollama-api}

Ollama exposes an OpenAI-compatible API at \texttt{http://localhost:11434},
enabling drop-in replacement for OpenAI API calls:

\begin{minted}{python}
from openai import OpenAI          # standard openai client

client = OpenAI(
    base_url="http://localhost:11434/v1",
    api_key="ollama",              # required but ignored by Ollama
)

response = client.chat.completions.create(
    model="mistral-finance",
    messages=[
        {"role": "system", "content": "You are a credit risk analyst."},
        {"role": "user",   "content": "Summarise the key covenants in: ..."},
    ],
    temperature=0.0,
)
print(response.choices[0].message.content)
\end{minted}

Existing Python code that calls the OpenAI API can be redirected to a local
Ollama instance by changing two lines---\texttt{base\_url} and
\texttt{api\_key}---with no other modifications.

% ---------------------------------------------------------------
\section{Production Inference Servers}
\label{app:hf:servers}

\begin{context}
Single-user local inference with Ollama or llama.cpp is sufficient for
individual research tasks.  When a model must serve multiple analysts
simultaneously---routing queries from a Slack bot, a web application, or an
automated pipeline---a production inference server is required.  The two
dominant open-source options are Hugging Face's Text Generation Inference
(TGI) and the community-developed vLLM.  They are architecturally distinct
and suit different workload profiles.
\end{context}

\subsection{Text Generation Inference (TGI)}
\label{app:hf:tgi}

TGI is Hugging Face's production-grade inference server and the backend for
Hugging Face Inference Endpoints.  Its design prioritises low tail latency in
interactive, single-user scenarios.  Key features:

\begin{itemize}
  \item Flash Attention 2 and Paged Attention for memory-efficient attention.
  \item Continuous batching, enabling multiple concurrent requests to share a
        single forward pass.
  \item Native support for AWQ, GPTQ, and bitsandbytes quantization.
  \item OpenAI-compatible API with streaming support.
\end{itemize}

\begin{minted}{bash}
docker run --gpus all \
  -p 8080:80 \
  -v $HOME/models:/data \
  ghcr.io/huggingface/text-generation-inference:latest \
  --model-id mistralai/Mistral-7B-Instruct-v0.3 \
  --quantize bitsandbytes-nf4 \
  --max-total-tokens 4096
\end{minted}

\subsection{vLLM and PagedAttention}
\label{app:hf:vllm}

vLLM \citep{kwon2023vllm} introduces \emph{PagedAttention}, an attention
mechanism inspired by virtual memory paging.  Rather than allocating a
contiguous block of GPU memory for each request's KV cache, PagedAttention
stores KV cache pages in non-contiguous blocks and maps them on demand.
This eliminates memory fragmentation, enables much larger effective batch
sizes, and delivers up to 24$\times$ higher throughput than TGI under
high-concurrency workloads.

\begin{minted}{bash}
pip install vllm

python -m vllm.entrypoints.openai.api_server \
    --model mistralai/Mistral-7B-Instruct-v0.3 \
    --quantization awq \
    --max-model-len 8192 \
    --port 8000
\end{minted}

In production, Stripe achieved a 73\% inference cost reduction by migrating
to vLLM, processing 50 million daily API calls on one-third of the original
GPU fleet \citep{kwon2023vllm}.

\subsection{Choosing an Inference Server}
\label{app:hf:server-choice}

\begin{table}[ht]
\centering
\caption{Inference server selection guide: TGI vs vLLM.}
\label{tab:server-choice}
\begin{tabular}{lll}
\toprule
\textbf{Criterion} & \textbf{TGI (Hugging Face)} & \textbf{vLLM} \\
\midrule
Throughput (high concurrency) & Good              & Excellent (up to 24$\times$ TGI) \\
Tail latency (single user)    & Excellent         & Good \\
Quantization support          & AWQ, GPTQ, BnB    & AWQ, GPTQ, FP8 \\
Kubernetes / Helm charts      & Official support  & Community Helm chart \\
HF Hub model compatibility    & Native            & Via \texttt{--model} flag \\
Best fit                      & Interactive apps  & Batch / high-throughput API \\
\bottomrule
\end{tabular}
\end{table}

For quantitative finance workflows, vLLM is typically the better choice when
running batch sentiment analysis, document extraction pipelines, or overnight
model monitoring jobs.  TGI is preferable for interactive analyst tools where
individual query latency matters more than aggregate throughput.

% ---------------------------------------------------------------
\section{Privacy, Compliance, and On-Premise Deployment}
\label{app:hf:privacy}

\subsection{Data Sovereignty and Regulatory Requirements}
\label{app:hf:sovereignty}

Cloud-hosted LLM APIs transmit prompt text to a third-party server.  For
financial institutions, this raises several regulatory concerns:

\begin{itemize}
  \item \textbf{Data residency} --- MiFID II, DORA, and national data
        protection laws may prohibit transferring client or trading data
        outside specified jurisdictions.
  \item \textbf{Model risk governance} --- SR 11-7 and equivalent frameworks
        require auditability of every model that influences a financial
        decision; a third-party API's internals are opaque.
  \item \textbf{Confidential information} --- Analyst reports, pending M\&A
        data, and portfolio holdings are material non-public information;
        submitting them to an external API may violate insider-trading
        regulations.
\end{itemize}

Local and on-premise deployment eliminates the data-transmission risk entirely:
all prompts and completions remain within the institution's controlled
environment.

\subsection{On-Premise Cost Analysis}
\label{app:hf:cost}

On-premise GPU servers have high upfront capital costs but dramatically lower
per-token compute costs for sustained workloads.  A server equipped with
8$\times$ NVIDIA H100 GPUs costs approximately \$0.87 per hour in electricity
and cooling; the equivalent cloud configuration costs approximately \$98.32 per
hour, a 113$\times$ difference \citep{openai2025codex}.  For a team running
continuous inference workloads, the breakeven point against cloud API pricing
is typically reached within 6--18 months of acquisition.

\begin{table}[ht]
\centering
\caption{On-premise vs cloud cost comparison for sustained LLM inference.}
\label{tab:onprem-cost}
\begin{tabular}{lrr}
\toprule
\textbf{Cost component} & \textbf{On-premise (8$\times$ H100)} & \textbf{Cloud (8$\times$ H100)} \\
\midrule
Compute (per hour)    & \$0.87  & \$98.32 \\
Capital amortisation  & \$3.50  & ---     \\
Effective total/hour  & \$4.37  & \$98.32 \\
Annual cost (24/7)    & \$38{,}281 & \$861{,}283 \\
\bottomrule
\end{tabular}
\end{table}

\subsection{Architecture Patterns for Finance}
\label{app:hf:arch}

Three deployment patterns are common in regulated finance environments:

\begin{description}
  \item[Air-gapped on-premise] All model weights, serving infrastructure, and
        data reside within the institution's data centre.  Maximum security;
        highest operational burden.  Required for most tier-1 trading system
        use cases.

  \item[Private cloud (VPC)] Model weights are loaded into a dedicated tenant
        environment on a major cloud provider's infrastructure.  Data remains
        within the institution's VPC; regulatory requirements for data residency
        are typically satisfied.

  \item[Hybrid] Public API for non-sensitive tasks (e.g., public document
        summarisation, general code assistance); on-premise or VPC deployment
        for tasks involving client or trading data.  Most cost-effective for
        institutions with a mix of sensitivity levels across workflows.
\end{description}

% ---------------------------------------------------------------
\section{Finance Research Workflows}
\label{app:hf:finance}

\subsection{Use Cases for Quants and Researchers}
\label{app:hf:use-cases}

\begin{description}
  \item[Earnings call sentiment] Run FinBERT or a fine-tuned Mistral variant
        over quarterly earnings transcripts to construct a sentiment factor.
        A Q4\_K\_M quantised model running on a workstation GPU can process
        500 transcripts per hour without any API cost.

  \item[Financial document extraction] Use a structured-output pipeline to
        extract numerical tables, covenant definitions, or risk factor summaries
        from PDF filings.  Local deployment eliminates concerns about disclosing
        confidential deal terms to a cloud provider.

  \item[RAG over internal research] Build a retrieval-augmented generation
        system over internal research notes using a local embedding model (e.g.,
        \texttt{BAAI/bge-m3}) and a local generation model.  Queries and
        responses never leave the institution's network.

  \item[Domain-adapted fine-tuning] Use QLoRA to adapt a 7B general model to
        the vocabulary and reasoning style of regulatory filings, credit
        agreements, or derivatives term sheets.  A 3-epoch fine-tune on 10{,}000
        examples completes in 4--6 hours on a single RTX 4090.

  \item[Stress testing and scenario generation] Prompt a local instruct model
        to generate hypothetical market narratives or macroeconomic scenarios
        for stress-testing model assumptions.
\end{description}

\begin{example}[End-to-end: local earnings-call sentiment factor]
\label{ex:hf-sentiment}
A researcher builds a monthly sentiment factor from quarterly earnings calls
using a fully local pipeline, with no data leaving the workstation.

\begin{minted}{python}
from transformers import pipeline
from datasets import Dataset
import pandas as pd

# 1. Load FinBERT once; all inference is local
clf = pipeline("text-classification",
               model="ProsusAI/finbert",
               device=0,           # GPU
               truncation=True, max_length=512)

# 2. Load earnings call transcripts (local CSV)
df = pd.read_csv("data/transcripts_2010_2024.csv")  # ticker, date, text

# 3. Batch inference (HF datasets handles batching automatically)
hf_ds = Dataset.from_pandas(df[["text"]])
results = hf_ds.map(
    lambda batch: {"label": [r["label"] for r in clf(batch["text"])],
                   "score": [r["score"]  for r in clf(batch["text"])]},
    batched=True, batch_size=32,
)

df["sentiment"] = results["label"]
df["confidence"] = results["score"]

# 4. Map label to numeric score, compute monthly averages
score_map = {"positive": 1, "neutral": 0, "negative": -1}
df["score_num"] = df["sentiment"].map(score_map)
factor = df.groupby(["ticker", pd.Grouper(key="date", freq="ME")])["score_num"].mean()

factor.to_csv("data/factors/finbert_sentiment.csv")
print(factor.describe())
\end{minted}

The entire pipeline runs on-premise, processes 500 transcripts per hour on
a single RTX 3080, and produces a factor ready for integration into a
Fama--French regression.
\end{example}

\subsection{Recommended Workflow Patterns}
\label{app:hf:patterns}

\begin{enumerate}
  \item \textbf{Start with the Hub, then localise.}  Prototype with the
        \texttt{pipeline()} API against the cloud Hub to identify the right
        model and prompt format.  Once the approach is validated, download
        the GGUF or safetensors weights and switch to local inference.

  \item \textbf{Match quantization to task sensitivity.}  Use Q4\_K\_M for
        classification and extraction tasks where 1\% quality loss is
        acceptable.  Use Q6\_K or F16 for tasks requiring precise numerical
        reasoning (e.g., extracting financial ratios from tables).

  \item \textbf{Fine-tune on domain data before deploying at scale.}  A
        QLoRA adapter trained on 5{,}000--10{,}000 finance-domain examples
        typically outperforms a much larger general model on domain-specific
        tasks at a fraction of the inference cost.

  \item \textbf{Use vLLM for batch, Ollama for interactive.}  Route overnight
        bulk pipelines through vLLM's OpenAI-compatible API; route analyst
        interactive queries through Ollama.  Both expose the same API schema,
        so switching requires only a \texttt{base\_url} change.

  \item \textbf{Enforce data-handling rules via the inference server.}  Set
        hard limits on \texttt{max\_tokens} and log all requests server-side.
        Never pass raw client data through a model server that is not subject
        to your institution's data-retention and access-control policies.

  \item \textbf{Version your models like code.}  Store model checkpoints and
        GGUF files in a content-addressed registry (e.g., DVC or MLflow
        artifacts).  Tag each deployment with the model hash, quantization
        level, and fine-tuning dataset version so that results are reproducible
        months later.
\end{enumerate}
